% PAGES: 148-154
% Continues after the c3 chunk (PDF pages 141-147 / folios 125-131).
% PDF page 147 (folio 131) ends section 4.6 "References" completely --
% the last reference paragraph ("...Lax and Zalcman [27].") is a complete
% sentence with no open environment, no open display, and no widow line
% carrying onto page 148. (The proof of the tridiagonal-matrix eigenvalue
% theorem that precedes "4.6 References" on that same page also closes
% cleanly with its QED box before the References section begins -- that
% proof and that whole section belong to the c3 chunk, not to this file.)
% PDF page 148 (folio 132) opens fresh with the "4.7 Exercises" section
% heading at the top of the page. No environment needs to be closed here;
% this file opens directly with \section{Exercises}.
%
% This is the whole of Chapter 4's exercise list (items 1-27, PDF pages
% 148-154 / folios 132-138), kept in a single \exercises environment so
% the shared enumerate counter is never reset mid-list -- this is also the
% last chunk of Chapter 4. PDF page 154 (folio 138) ends with the hint to
% exercise 27; the remainder of that page is blank (end of chapter). No
% environment is left open at the end of this file.

\section{Exercises}

\begin{exercises}

\item Let $A$ and $B$ be square matrices of the same size. Show that, if $v$ is
an eigenvector of both $A$ and $B$, then $v$ is also an eigenvector of the
matrix $M = c_1 A + c_2 B$, where $c_1$ and $c_2$ are constants. What is the
eigenvalue of $M$ corresponding to the eigenvector $v$?

\item Let $\lambda$ and $v$ be an eigenvalue and the corresponding
eigenvector of the matrix $A$. Let $d$ be a constant number, and let $I$ be
the identity matrix.
\begin{subparts}
\item Show that $v$ is an eigenvector of the matrix $B = dI + A + A^2$, and
find the corresponding eigenvalue.
\item If $A$ is a nonsingular matrix, show that $v$ is an eigenvector of the
matrix $M = dI + A + A^{-1}$, and find the corresponding eigenvalue.
\end{subparts}

\item Let $A$ be a $3 \times 3$ nonsingular matrix with eigenvalues $-2$, $1$,
and $3$. What are the eigenvalues of $A^2 + 2I - 3A^{-1}$?

\item
\begin{subparts}
\item Show that the eigenvalues of the matrix
\[
A \;=\; \begin{pmatrix} 8 & -18 & -30 & -24\\ 18 & -37 & -60 & -48\\ -9 & 18 & 29 & 24\\ 0 & 0 & 0 & -1 \end{pmatrix}
\]
are $-1$, with multiplicity 3, and 2, with multiplicity 1.

Show that $\begin{pmatrix} 0\\ 1\\ 1\\ -2 \end{pmatrix}$,
$\begin{pmatrix} 2\\ 2\\ -3\\ 3 \end{pmatrix}$,
$\begin{pmatrix} 2\\ 1\\ 0\\ 0 \end{pmatrix}$ are three linearly independent
eigenvectors corresponding to the eigenvalue $-1$, and show that there exists
only one linearly independent eigenvector of the matrix $A$ corresponding to
the eigenvalue 2, e.g., $\begin{pmatrix} 1\\ 2\\ -1\\ 0 \end{pmatrix}$.

\item Show that the eigenvalues of the matrix
\[
B \;=\; \begin{pmatrix} 10 & -20 & -32 & -26\\ 18 & -41 & -68 & -54\\ -14 & 19 & 26 & 23\\ 7 & 1 & 9 & 4 \end{pmatrix}
\]
are $-1$, with multiplicity 3, and 2, with multiplicity 1.

Show that there exists only one linearly independent eigenvector of the
matrix $B$ corresponding to the eigenvalue $-1$, e.g.,
$\begin{pmatrix} 0\\ -1\\ -1\\ 2 \end{pmatrix}$, and show that
$\begin{pmatrix} 1\\ 2\\ -1\\ 0 \end{pmatrix}$ is an eigenvector corresponding
to the eigenvalue 2.
\end{subparts}

\item Let $A = \begin{pmatrix} 2 & -1\\ 1 & 4 \end{pmatrix}$.
\begin{subparts}
\item What is the characteristic polynomial of $A$?
\item Find the eigenvalues and the eigenvectors of $A$.
\end{subparts}

\item
\begin{subparts}
\item Find the eigenvalues and the eigenvectors of the lower triangular
matrix
\[
L \;=\; \begin{pmatrix} 2 & 0 & 0\\ 1 & -3 & 0\\ -1 & 2 & -1 \end{pmatrix}.
\]
\item Find the eigenvalues and the eigenvectors of the upper triangular
matrix
\[
U \;=\; \begin{pmatrix} -2 & -1 & 3\\ 0 & 1 & 2\\ 0 & 0 & 3 \end{pmatrix}.
\]
\end{subparts}

\item Let $A$ be a square matrix such that $A^2 = A$. Show that any
eigenvalue of $A$ is either 0 or 1.

\noindent Note: A matrix $A$ with the property that $A^2 = A$ is called an
idempotent matrix.

\item Let $A$ be a square matrix with the property that there exists a
positive integer $p$ such that $A^p = 0$. Show that any eigenvalue of $A$
must be equal to 0.

\noindent Note: A matrix $A$ with the property that $A^p = 0$ for a positive
integer $p$ is called a nilpotent matrix.

\item Let $v \neq 0$ be a column vector of size $n$, and let $A = vv^t$ be an
$n \times n$ matrix.
\begin{subparts}
\item Show that the matrix $A$ has exactly one non-zero eigenvalue.
\item Find the eigenvalues and the eigenvectors of $A$.
\end{subparts}

\item Find the eigenvalues and the eigenvectors of the $n \times n$ matrix
\[
\begin{pmatrix} d & 1 & \dots & 1\\ 1 & d & \dots & 1\\ \vdots & \vdots & \ddots & \vdots\\ 1 & 1 & \dots & d \end{pmatrix},
\]
where $d \in \R$ is a constant.

\noindent Hint: Note that
\begin{align*}
\begin{pmatrix} d & 1 & \dots & 1\\ 1 & d & \dots & 1\\ \vdots & \vdots & \ddots & \vdots\\ 1 & 1 & \dots & d \end{pmatrix}
&= (d-1)I \;+\; \begin{pmatrix} 1 & 1 & \dots & 1\\ 1 & 1 & \dots & 1\\ \vdots & \vdots & \ddots & \vdots\\ 1 & 1 & \dots & 1 \end{pmatrix}\\
&= (d-1)I \;+\; \bfone \cdot \bfone^t,
\end{align*}
where $I$ is the identity matrix and $\bfone = \begin{pmatrix} 1\\ 1\\ \vdots\\ 1 \end{pmatrix}$.

\item
\begin{subparts}
\item Let $A$ be a symmetric matrix of size $n$, and let $\lambda_i$, $i =
1:n$, be the eigenvalues of $A$, with corresponding eigenvectors $v_i$, $i =
1:n$. What are the eigenvalues and the eigenvectors of $A^2$?
\item Let $A$ be a symmetric matrix of size $n$. Let $\phi_i$, $i = 1:n$, be
the eigenvalues of the matrix $A^2$, with corresponding eigenvectors $w_i$, $i
= 1:n$. What can you say about the eigenvalues and the eigenvectors of $A$?
\end{subparts}

\item Let $A$ be a square matrix of size $n$, and let $S$ be a nonsingular
matrix of size $n$.
\begin{subparts}
\item Let $\lambda$ and $v$ be an eigenvalue and the corresponding
eigenvector of $A$. Show that $\lambda$ is also an eigenvalue of the matrix
$S^{-1}AS$. What is the corresponding eigenvector?
\item Show that the matrix $S^{-1}AS$ has the same characteristic polynomial
as $A$, i.e., show that
\[
P_{S^{-1}AS}(t) \;=\; P_A(t), \quad \forall\, t \in \R.
\]
\end{subparts}

\item Let $A$ be a square matrix with real entries. If $\lambda = a + ib$ is a
complex eigenvalue of $A$ (i.e., with $b \neq 0$), show that $\overline{\lambda}
= a - ib$, the complex conjugate of $\lambda$, is also an eigenvalue of $A$.

\item The $2 \times 2$ matrix $A$ has eigenvalues 1 and $-2$ with
corresponding eigenvectors $\begin{pmatrix} -1\\ 2 \end{pmatrix}$ and
$\begin{pmatrix} 3\\ 1 \end{pmatrix}$. If $v = \begin{pmatrix} 2\\ -3
\end{pmatrix}$, find $Av$.

\item Let $A = \begin{pmatrix} -1 & 2\\ 2 & 2 \end{pmatrix}$.
\begin{subparts}
\item Find the eigenvalues and the eigenvectors of the matrix $A$.
\item What is the diagonal form of $A$?
\item Compute $A^{12}$.
\end{subparts}

\item Let $A = \begin{pmatrix} a & b\\ c & d \end{pmatrix}$ be a $2 \times 2$
matrix, and let
\[
P_A(t) \;=\; t^2 - (a+d)t + (ad-bc)
\]
be the characteristic polynomial associated to $A$.

Show that $P_A(A) = 0$, i.e., show that
\[
A^2 - (a+d)A + (ad-bc)I \;=\; 0.
\]

\noindent Note: This is the $2 \times 2$ case of the Cayley--Hamilton
theorem which states that $P_A(A) = 0$ for any square matrix $A$.

\item Let $A$ and $B$ be square matrices of the same size.
\begin{subparts}
\item If $A$ is nonsingular, show that the matrices $AB$ and $BA$ have the
same characteristic polynomial, i.e., show that
\[
\det(tI - AB) \;=\; \det(tI - BA).
\]
\item Show that the characteristic polynomial of $AB$ is the same as the
characteristic polynomial of $BA$ even if $A$ is a singular matrix.

\noindent Hint: If $A$ is a singular matrix, it is possible to find a number
$\epsilon > 0$ as small as needed such that the matrix $A + \epsilon I$ is
nonsingular.
\item Show that the matrices $AB$ and $BA$ have the same eigenvalues.
\end{subparts}

\item
\begin{subparts}
\item Let $A$ and $B$ be square matrices of the same size. Show that the
traces of the matrices $AB$ and $BA$ are equal, i.e., show that $\tr(AB) =
\tr(BA)$.
\item Show that you cannot find two $n \times n$ matrices $A$ and $B$ such
that
\[
AB - BA \;=\; I,
\]
where $I$ is the $n \times n$ identity matrix.
\end{subparts}

\item Recall that the matrix
\[
B_4 \;=\; \begin{pmatrix} 2 & -1 & 0 & 0\\ -1 & 2 & -1 & 0\\ 0 & -1 & 2 & -1\\ 0 & 0 & -1 & 2 \end{pmatrix}
\]
has four eigenvectors $v_1$, $v_2$, $v_3$, $v_4$ given by
\[
v_j(i) \;=\; \sin\left(\frac{ij\pi}{5}\right), \quad \forall\, i = 1:4,
\]
for $j = 1:4$.

Are the vectors $v_1$, $v_2$, $v_3$, and $v_4$ orthogonal and of norm 1?

\item Show that the matrix $\begin{pmatrix} 3 & -1 & -1\\ 0 & 2 & 2\\ 0 & 1 &
1 \end{pmatrix}$ is a weakly diagonally dominant singular matrix.

\item Show that the matrix
\[
\begin{pmatrix} -5 & -1 & 0.25 & -1\\ 2 & 4 & -1 & 3\\ 0.5 & -2 & 3 & -1\\ 1 & 0.5 & 0 & -6 \end{pmatrix}
\]
is nonsingular.

\item Let $A$ be an $n \times n$ matrix. For every $j = 1:n$, let
\[
R_j \;=\; \sum_{k=1:n,\, k \neq j} |A(j,k)|,
\]
and denote by $D_j$ the disc of center $A(j,j)$ and radius $R_j$, i.e.,
\[
D_j \;=\; \set{z \in \C \ \text{such that} \ |z - A(j,j)| \leq R_j};
\]
note that $D_j$ is called a Gershgorin disk corresponding to the matrix $A$.

A more general form of Gershgorin's theorem states that, if a Gershgorin disk
$D_i$ is disjoint from the union of the other $n-1$ Gershgorin disks of $A$,
then exactly one eigenvalue of the matrix $A$ is in the disk $D_i$.

Use this result to show that all the eigenvalues of the matrix
\[
\begin{pmatrix} -4 & 1 & 0 & -0.5\\ 0 & 0.1 & 0 & -0.2\\ 1 & 2 & 5 & -1\\ 0.25 & -0.15 & 0.1 & -1 \end{pmatrix}
\]
are real numbers.

\item Show that all the eigenvalues of the matrix
\[
\begin{pmatrix} 2 & 0.0012 & -0.0003 & 0.0015\\ -0.0002 & -1.25 & 0.0010 & -0.0001\\ 0 & 0.0016 & 3 & 0.0009\\ -0.0011 & -0.0008 & -0.0002 & -2.5 \end{pmatrix}
\]
are real numbers, and find estimates for the values of the eigenvalues of the
matrix with 0.005 accuracy.

\item Let $A$ be an $n \times n$ matrix given by
\begin{align*}
A(i,i) &= 2, \quad \forall\, i = 1:n;\\
A(i,i-1) &= 1, \quad \forall\, i = 2:n;\\
A(j,k) &= 0, \quad \text{otherwise}.
\end{align*}
Find the eigenvalues and the eigenvectors of $A$.

\item Let $J$ be an $n \times n$ matrix given by
\begin{align*}
A(i,i) &= a, \quad \forall\, i = 1:n;\\
A(i,i+1) &= b, \quad \forall\, i = 1:(n-1);\\
A(j,k) &= 0, \quad \text{otherwise},
\end{align*}
where $a, b \in \R$ are constants. Find the eigenvalues and the eigenvectors
of $J$.

\noindent Note: The matrix $J$ is called a Jordan block if $b = 1$.

\item The 2--norm of a symmetric matrix $A$ is given by
\begin{equation}
\norm{A}_2 \;=\; \max_{\lambda \text{ eigenvalue of } A} |\lambda|,
\end{equation}
i.e., $\norm{A}_2$ is equal to the largest absolute value of all the
eigenvalues of $A$.

The Forward Euler finite difference discretization of the heat PDE is
convergent if and only if $\norm{A_N}_2 \leq 1$ for all $N \geq 2$, where
$A_N$ is the $N \times N$ matrix given by
\[
A_N \;=\; \begin{pmatrix} 1-2\alpha & -\alpha & \dots & 0\\ -\alpha & \ddots & \ddots & \vdots\\ \vdots & \ddots & \ddots & -\alpha\\ 0 & \dots & -\alpha & 1-2\alpha \end{pmatrix},
\]
with $\alpha > 0$ a positive constant.

Show that $\norm{A_N}_2 \leq 1$ for all $N \geq 2$ if and only if $0 < \alpha
\leq \frac{1}{2}$.

\noindent Hint: Recall from Lemma 4.11 that the matrix
\[
T_N \;=\; \begin{pmatrix} d & -a & \dots & 0\\ -a & \ddots & \ddots & \vdots\\ \vdots & \ddots & \ddots & -a\\ 0 & \dots & -a & d \end{pmatrix}
\]
has the following $N$ different eigenvalues:
\[
\lambda_j \;=\; d - 2a\cos\left(\frac{\pi j}{N+1}\right), \quad \text{for} \ j = 1:N.
\]

\item The Backward Euler finite difference discretization of the heat PDE is
convergent if and only if $\norm{A_N^{-1}}_2 \leq 1$ for all $N \geq 2$, where
\[
A_N \;=\; \begin{pmatrix} 1+2\alpha & \alpha & \dots & 0\\ \alpha & \ddots & \ddots & \vdots\\ \vdots & \ddots & \ddots & \alpha\\ 0 & \dots & \alpha & 1+2\alpha \end{pmatrix},
\]
with $\alpha > 0$ a positive constant.

Show that $\norm{A_N^{-1}}_2 \leq 1$ for all $N \geq 2$ for any $\alpha > 0$.

\noindent Hint: Use (4.54) to show that, if $A$ is a symmetric matrix, then
\[
\norm{A^{-1}}_2 \;=\; 1/\bigl(\min_{\lambda \text{ eigenvalue of } A} |\lambda|\bigr).
\]

\end{exercises}

% PDF page 154 (folio 138) ends here -- the hint to exercise 27 is the last
% content on the page; the remainder of the page is blank. This is the end
% of Chapter 4 "Eigenvalues and Eigenvectors." No environment is left open.
